\chapter{Théorème de Thalès}

\noindent\textit{Le théorème de Thalès relie les longueurs dans une configuration de droites
parallèles. Il permet de calculer une longueur inaccessible (hauteur d'un arbre, largeur
d'une rivière), de démontrer qu'un parallélisme existe, et de comprendre l'effet d'un
agrandissement ou d'une réduction sur les périmètres et les aires.}
\vspace{8pt}

\connecteBanner

\section{Configuration de Thalès}

\begin{definition}[Configuration de Thalès]
Dans un triangle $ABC$, on place un point $D$ sur $[AB]$ et un point $E$ sur $[AC]$. On dit
que $D$ et $E$ forment une \textbf{configuration de Thalès} avec $B$ et $C$ lorsque
$(DE)\parallel(BC)$ : on obtient deux triangles « emboîtés » $ADE$ et $ABC$.
\end{definition}

\begin{illus}
\begin{tikzpicture}[scale=0.9]
\coordinate (A) at (0,0);
\coordinate (B) at (7,0);
\coordinate (C) at (2,5);
\coordinate (D) at (4.67,0);
\coordinate (E) at (1.33,3.33);
\draw[onbo,line width=1.1pt] (A) -- (B) -- (C) -- cycle;
\draw[onbblue,line width=1.3pt] (D) -- (E);
\fill (A) circle(1.6pt) (B) circle(1.6pt) (C) circle(1.6pt) (D) circle(1.6pt) (E) circle(1.6pt);
\node[below left] at (A) {$A$};
\node[below right] at (B) {$B$};
\node[above] at (C) {$C$};
\node[below] at (D) {$D$};
\node[above left] at (E) {$E$};
\node[below] at ($(A)!0.5!(D)$) {\small $4$};
\node[below] at ($(D)!0.5!(B)$) {\small $2$};
\node[left] at ($(A)!0.5!(E)$) {\small $6$};
\node[right] at ($(E)!0.5!(C)$) {\small $3$};
\node[onbblue,above] at ($(D)!0.5!(E)$) {\small $5$};
\node[onbo,below right] at ($(B)!0.5!(C)$) {\small $7{,}5$};
\end{tikzpicture}
\fig{Figure — Configuration de Thalès : $(DE)\parallel(BC)$ dans le triangle $ABC$.}
\end{illus}

\section{Théorème de Thalès}

\begin{propriete}[Théorème de Thalès]
Si $D\in[AB]$, $E\in[AC]$ et $(DE)\parallel(BC)$, alors
\[
\frac{AD}{AB}=\frac{AE}{AC}=\frac{DE}{BC}.
\]
\end{propriete}

\begin{illus}
\begin{tikzpicture}[>=Stealth,every node/.style={font=\small}]
\node[draw=onbo,fill=onbsoft,rounded corners=3pt,minimum width=3.6cm,minimum height=0.7cm] (a1) {$(DE)\parallel(BC)$};
\node[draw=onbblue,fill=white,rounded corners=3pt,minimum width=5.4cm,minimum height=0.9cm,below=8mm of a1] (a2) {$\dfrac{AD}{AB}=\dfrac{AE}{AC}=\dfrac{DE}{BC}$};
\draw[->,onbmuted,line width=1pt] (a1) -- (a2);
\end{tikzpicture}
\fig{Figure — Le théorème de Thalès : trois rapports égaux dans le triangle emboîté.}
\end{illus}

\begin{exemple}
Avec $AD=4$, $AB=6$, $AC=9$, $BC=7{,}5$ : le rapport est $\dfrac{AD}{AB}=\dfrac46=\dfrac23$,
donc $AE=\dfrac23\times9=6$ et $DE=\dfrac23\times7{,}5=5$.
\end{exemple}

\section{Réciproque du théorème de Thalès}

\begin{propriete}[Réciproque de Thalès]
Si $D\in[AB]$, $E\in[AC]$ (avec $D,E$ du même côté de $A$) et si
$\dfrac{AD}{AB}=\dfrac{AE}{AC}$ (soit $AD\times AC=AB\times AE$), alors $(DE)\parallel(BC)$.
\end{propriete}

\begin{illus}
\begin{tikzpicture}[scale=0.9]
\coordinate (A) at (0,0);
\coordinate (B) at (7,0);
\coordinate (C) at (2,5);
\coordinate (D) at (2.8,0);
\coordinate (E) at (0.75,1.875);
\draw[onbo,line width=1.1pt] (A) -- (B) -- (C) -- cycle;
\draw[onbblue,dashed,line width=1.3pt] (D) -- (E);
\fill (A) circle(1.6pt) (B) circle(1.6pt) (C) circle(1.6pt) (D) circle(1.6pt) (E) circle(1.6pt);
\node[below left] at (A) {$A$};
\node[below right] at (B) {$B$};
\node[above] at (C) {$C$};
\node[below] at (D) {$D$};
\node[left] at (E) {$E$};
\node[below] at ($(A)!0.5!(D)$) {\small $2$};
\node[left] at ($(A)!0.5!(E)$) {\small $3$};
\end{tikzpicture}
\fig{Figure — $AD=2$, $AB=5$, $AE=3$, $AC=8$ : $(DE)$ et $(BC)$ ne sont pas parallèles.}
\end{illus}

\begin{exemple}
$AD=3$, $AB=5$, $AE=6$, $AC=10$ : $AD\times AC=3\times10=30$ et $AB\times AE=5\times6=30$.
Les produits sont égaux, donc $(DE)\parallel(BC)$.
\end{exemple}

\begin{remarque}
Comparer les \textbf{produits en croix} $AD\times AC$ et $AB\times AE$ évite les décimales.
Si les produits sont \textbf{différents}, les droites ne sont \textbf{pas} parallèles (c'est
la contraposée de la réciproque).
\end{remarque}

\section{Configuration « papillon »}

\begin{propriete}[Droites sécantes]
Si deux droites $(AC)$ et $(BD)$ se coupent en $O$ (avec $A,O,C$ et $B,O,D$ alignés dans cet
ordre croisé) et si $(AB)\parallel(CD)$, alors
\[
\frac{OA}{OC}=\frac{OB}{OD}=\frac{AB}{CD}.
\]
\end{propriete}

\begin{exemple}
$OA=4$, $OB=6$, $OC=10$, $(AB)\parallel(CD)$ : $\dfrac{OA}{OC}=\dfrac{OB}{OD}$, donc
$OD=\dfrac{OB\times OC}{OA}=\dfrac{6\times10}4=15$.
\end{exemple}

\section{Agrandissement, réduction : effet sur périmètres et aires}

\begin{propriete}[Coefficient d'agrandissement/réduction]
Lorsqu'une figure est agrandie ou réduite avec un coefficient $k$ ($k>1$ : agrandissement ;
$0<k<1$ : réduction), toutes les longueurs sont multipliées par $k$, le \textbf{périmètre}
est multiplié par $k$, et l'\textbf{aire} est multipliée par $k^{2}$.
\end{propriete}

\begin{illus}
\begin{tikzpicture}
\draw[fill=onbsoft,draw=onbo,line width=1pt] (0,0) rectangle (3,2);
\node at (1.5,1) {\small $k=1$};
\draw[fill=onbsoft,draw=onbblue,line width=1pt] (5,0) rectangle (9.5,3);
\node at (7.25,1.5) {\small $k=1{,}5$};
\node[below] at (1.5,-0.3) {\small périmètre $10$, aire $6$};
\node[below] at (7.25,-0.3) {\small périmètre $15$, aire $13{,}5$};
\end{tikzpicture}
\fig{Figure — Un agrandissement $k=1{,}5$ multiplie le périmètre par $1{,}5$ et l'aire par
$1{,}5^{2}=2{,}25$.}
\end{illus}

\begin{exemple}
Un rectangle de périmètre $10$ et d'aire $6$ est agrandi avec $k=1{,}5$ : nouveau périmètre
$=10\times1{,}5=15$, nouvelle aire $=6\times1{,}5^{2}=6\times2{,}25=13{,}5$.
\end{exemple}

% ═══════════════════════════════════════════════════════════════
\section{L'essentiel à retenir}

\begin{synthese}
\textbf{Configuration.} $D\in[AB]$, $E\in[AC]$, $(DE)\parallel(BC)$ : triangle $ADE$
emboîté dans $ABC$.\\[4pt]
\textbf{Théorème de Thalès.} $(DE)\parallel(BC)\Rightarrow\dfrac{AD}{AB}=\dfrac{AE}{AC}=\dfrac{DE}{BC}$.\\[4pt]
\textbf{Réciproque.} $AD\times AC=AB\times AE\Rightarrow(DE)\parallel(BC)$ (produits en
croix égaux).\\[4pt]
\textbf{Papillon.} Droites sécantes en $O$, $(AB)\parallel(CD)$ :
$\dfrac{OA}{OC}=\dfrac{OB}{OD}=\dfrac{AB}{CD}$.\\[4pt]
\textbf{Agrandissement/réduction.} Coefficient $k$ : longueurs $\times k$, périmètre
$\times k$, aire $\times k^{2}$.\\[4pt]
\textbf{Piège fréquent.} N'utiliser le théorème que si le parallélisme est \textbf{donné}
ou déjà \textbf{prouvé} ; pour prouver un parallélisme, c'est la réciproque qu'il faut
utiliser (comparer les produits en croix, pas les rapports décimaux).
\end{synthese}

% ═══════════════════════════════════════════════════════════════
\section{Méthodes \& savoir-faire}

\methode{Identifier une configuration de Thalès et écrire les rapports égaux}
Repérer le sommet commun $A$, les points $D,E$ alignés avec $B,C$, vérifier que
$(DE)\parallel(BC)$ est donné, puis écrire $\dfrac{AD}{AB}=\dfrac{AE}{AC}=\dfrac{DE}{BC}$.
\begin{exemple}
Triangle $ABC$, $D\in[AB]$, $E\in[AC]$, $(DE)\parallel(BC)$ : on écrit directement
$\dfrac{AD}{AB}=\dfrac{AE}{AC}=\dfrac{DE}{BC}$.
\end{exemple}

\methode{Calculer une longueur inconnue avec le théorème de Thalès}
Calculer le rapport connu, puis l'appliquer à la longueur cherchée.
\begin{exemple}
$AD=6$, $AB=9$, $AC=15$ : rapport $\dfrac69=\dfrac23$, donc $AE=\dfrac23\times15=10$.
\end{exemple}

\methode{Utiliser la réciproque pour prouver un parallélisme}
Calculer les produits en croix $AD\times AC$ et $AB\times AE$ : s'ils sont égaux, les droites
sont parallèles.
\begin{exemple}
$AD=4$, $AB=6$, $AE=6$, $AC=9$ : $4\times9=36$ et $6\times6=36$. Égaux, donc
$(DE)\parallel(BC)$.
\end{exemple}

\methode{Utiliser la réciproque pour prouver que deux droites ne sont pas parallèles}
Calculer les mêmes produits en croix : s'ils sont différents, les droites ne sont pas
parallèles.
\begin{exemple}
$AD=5$, $AB=8$, $AE=6$, $AC=10$ : $5\times10=50$ et $8\times6=48$. Différents, donc $(DE)$
et $(BC)$ ne sont pas parallèles.
\end{exemple}

\methode{Résoudre un problème de configuration « papillon »}
Repérer le point d'intersection $O$, identifier les deux triangles opposés par le sommet,
puis écrire $\dfrac{OA}{OC}=\dfrac{OB}{OD}$.
\begin{exemple}
$OA=5$, $OC=15$, $OB=8$ : $OD=\dfrac{OB\times OC}{OA}=\dfrac{8\times15}5=24$.
\end{exemple}

\methode{Calculer l'effet d'un agrandissement/réduction sur un périmètre}
Multiplier le périmètre initial par le coefficient $k$.
\begin{exemple}
Périmètre $18$, coefficient $k=\dfrac23$ : nouveau périmètre $=18\times\dfrac23=12$.
\end{exemple}

\methode{Calculer l'effet d'un agrandissement/réduction sur une aire}
Multiplier l'aire initiale par $k^{2}$.
\begin{exemple}
Aire $36$, coefficient $k=\dfrac23$ : nouvelle aire $=36\times\left(\dfrac23\right)^{2}=36\times\dfrac49=16$.
\end{exemple}

% ═══════════════════════════════════════════════════════════════
\section{Exercices d'application}
\begin{multicols}{2}

\rubrique{Identifier et calculer avec le théorème de Thalès}
\exo{}\diff{1} Triangle $ABC$, $D\in[AB]$, $E\in[AC]$, $(DE)\parallel(BC)$. $AD=6$,
$AB=9$, $AC=15$, $BC=12$. Calculer $AE$ et $DE$.

\exo{}\diff{2} Mêmes conditions. $AD=3$, $AB=4$, $AC=12$, $BC=16$. Calculer $AE$ et $DE$.

\exo{}\diff{2} Mêmes conditions. $AD=1$, $AB=3$, $AC=9$, $BC=12$. Calculer $AE$ et $DE$.

\exo{}\diff{2} Mêmes conditions. $AD=4$, $AB=5$, $AC=10$, $BC=15$. Calculer $AE$ et $DE$.

\rubrique{Réciproque : prouver un parallélisme}
\exo{}\diff{2} $D\in[AB]$, $E\in[AC]$ avec $AD=3$, $AB=5$, $AE=6$, $AC=10$. Les droites
$(DE)$ et $(BC)$ sont-elles parallèles ?

\exo{}\diff{2} $AD=4$, $AB=6$, $AE=6$, $AC=9$. Les droites $(DE)$ et $(BC)$ sont-elles
parallèles ?

\exo{}\diff{2} $AD=2$, $AB=5$, $AE=3$, $AC=8$. Les droites $(DE)$ et $(BC)$ sont-elles
parallèles ?

\exo{}\diff{2} $AD=5$, $AB=8$, $AE=6$, $AC=10$. Les droites $(DE)$ et $(BC)$ sont-elles
parallèles ?

\rubrique{Papillon et agrandissement/réduction}
\exo{}\diff{2} Les droites $(AC)$ et $(BD)$ se coupent en $O$, avec $(AB)\parallel(CD)$.
$OA=4$, $OB=6$, $OC=10$. Calculer $OD$.

\exo{}\diff{1} Un rectangle de côtés $4$ cm et $6$ cm est agrandi avec le coefficient
$k=1{,}5$. Calculer le nouveau périmètre et la nouvelle aire.

\exo{}\diff{2} Un triangle a un périmètre de $18$ cm et une aire de $36\text{ cm}^2$. Il
est réduit avec le coefficient $k=\dfrac23$. Calculer le nouveau périmètre et la nouvelle
aire.

\exo{}\diff{1} Un terrain de périmètre $80$ m est agrandi avec le coefficient $k=1{,}25$.
Calculer le nouveau périmètre.

\rubrique{Problèmes concrets}
\exo{}\diff{2} Pour mesurer la hauteur d'un arbre, un élève plante un bâton vertical de
$1{,}2$ m dont l'ombre mesure $1{,}8$ m. Au même moment, l'ombre de l'arbre mesure $10{,}5$
m. Calculer la hauteur de l'arbre.

\exo{}\diff{2} Pour estimer la largeur $BC$ d'une rivière, un géomètre relève, dans une
configuration de Thalès, $AD=6$ m, $AB=15$ m, $AC=20$ m et $DE=8$ m, avec $(DE)\parallel(BC)$.
Calculer $BC$ et $AE$.

\exo{}\diff{2} Sur le plan d'une maison, une pièce a une aire de $12\text{ cm}^2$. Le plan
est à l'échelle $1/100$ (coefficient d'agrandissement réel $100$). Calculer l'aire réelle de
la pièce, en $\text{m}^2$.

\exo{}\diff{2} Deux allées d'un jardin se croisent en un point $O$, avec $(AB)\parallel(CD)$.
$OA=5$, $OB=8$, $OD=24$. Calculer $OC$.

% ═══════════════════════════════════════════════════════════════
\end{multicols}

\section{Exercices d'entraînement}
\begin{multicols}{2}

\rubrique{Calculs de longueurs}
\exo{}\diff{2} $AD=2$, $AB=3$, $AC=15$, $BC=9$, $(DE)\parallel(BC)$. Calculer $AE$ et $DE$.

\exo{}\diff{2} $AD=3$, $AB=5$, $AC=20$, $BC=15$, $(DE)\parallel(BC)$. Calculer $AE$ et $DE$.

\exo{}\diff{2} $AD=2$, $AB=7$, $AC=21$, $BC=14$, $(DE)\parallel(BC)$. Calculer $AE$ et $DE$.

\exo{}\diff{2} $AD=5$, $AB=7$, $AC=14$, $BC=21$, $(DE)\parallel(BC)$. Calculer $AE$ et $DE$.

\exo{}\diff{2} $AD=1$, $AB=4$, $AC=16$, $BC=20$, $(DE)\parallel(BC)$. Calculer $AE$ et $DE$.

\rubrique{Réciproque : parallèle ou non}
\exo{}\diff{2} $AD=4$, $AB=7$, $AE=8$, $AC=14$. $(DE)$ et $(BC)$ sont-elles parallèles ?

\exo{}\diff{2} $AD=3$, $AB=8$, $AE=5$, $AC=12$. $(DE)$ et $(BC)$ sont-elles parallèles ?

\exo{}\diff{2} $AD=6$, $AB=10$, $AE=9$, $AC=15$. $(DE)$ et $(BC)$ sont-elles parallèles ?

\exo{}\diff{2} $AD=2$, $AB=9$, $AE=5$, $AC=20$. $(DE)$ et $(BC)$ sont-elles parallèles ?

\rubrique{Papillon}
\exo{}\diff{2} $(AC)$ et $(BD)$ se coupent en $O$, $(AB)\parallel(CD)$. $OA=3$, $OB=11$,
$OC=9$. Calculer $OD$.

\exo{}\diff{2} $(AC)$ et $(BD)$ se coupent en $O$, $(AB)\parallel(CD)$. $OA=6$, $OB=9$,
$OC=8$. Calculer $OD$.

\rubrique{Agrandissement, réduction}
\exo{}\diff{1} Une figure de périmètre $16$ cm et d'aire $20\text{ cm}^2$ est agrandie
avec le coefficient $k=2$. Calculer le nouveau périmètre et la nouvelle aire.

\exo{}\diff{2} Une figure de périmètre $24$ cm et d'aire $32\text{ cm}^2$ est réduite avec
le coefficient $k=0{,}5$. Calculer le nouveau périmètre et la nouvelle aire.

\exo{}\diff{2} Une figure de périmètre $21$ cm et d'aire $27\text{ cm}^2$ est agrandie avec
le coefficient $k=\dfrac43$. Calculer le nouveau périmètre et la nouvelle aire.

\exo{}\diff{2} Une figure de périmètre $28$ cm et d'aire $40\text{ cm}^2$ est agrandie avec
le coefficient $k=\dfrac54$. Calculer le nouveau périmètre et la nouvelle aire.

\rubrique{Problèmes concrets}
\exo{}\diff{2} Un bâton de $1{,}5$ m planté verticalement projette une ombre de $2$ m. Au
même moment, l'ombre d'un poteau électrique mesure $8$ m. Calculer la hauteur du poteau.

\exo{}\diff{2} Pour estimer la largeur d'une rivière, un géomètre relève $AD=8$ m, $AB=20$
m, $AC=25$ m et $DE=10$ m, dans une configuration de Thalès avec $(DE)\parallel(BC)$.
Calculer la largeur $BC$ de la rivière.

\exo{}\diff{2} Sur un plan à l'échelle $1/50$, un jardin carré a une aire de $100\text{ cm}^2$.
Calculer l'aire réelle du jardin, en $\text{m}^2$ (coefficient d'agrandissement réel $50$).

\exo{}\diff{2} Un jardin carré de $25\text{ m}^2$ est représenté sur un plan à l'échelle
$1/50$. Calculer l'aire du jardin sur le plan, en $\text{cm}^2$.

% ═══════════════════════════════════════════════════════════════
\end{multicols}

\section{Sujets type BEPC}

\sujetbac[Thalès, réciproque et agrandissement d'un terrain]
\begin{enumerate}[label=\arabic*)]
\item Triangle $ABC$, $D\in[AB]$, $E\in[AC]$, $(DE)\parallel(BC)$. $AD=8$, $AB=12$,
$AC=18$, $BC=15$. Calculer $AE$ et $DE$.
\item Dans un triangle $MNP$, $R\in[MN]$, $S\in[MP]$, avec $MR=5$, $MN=8$, $MS=9$, $MP=15$.
Les droites $(RS)$ et $(NP)$ sont-elles parallèles ? Justifier par un calcul de produits en
croix.
\item Un terrain a un périmètre de $40$ m et une aire de $75\text{ m}^2$. On l'agrandit
avec un coefficient $k=1{,}2$.
\begin{enumerate}[label=\alph*)]
\item Calculer le nouveau périmètre.
\item Calculer la nouvelle aire.
\end{enumerate}
\end{enumerate}

\sujetbac[Thalès et mesures de terrain — Ebolowa]
\begin{enumerate}[label=\arabic*)]
\item Triangle $ABC$, $D\in[AB]$, $E\in[AC]$, $(DE)\parallel(BC)$. $AD=5$, $AB=7$,
$AC=14$, $BC=21$. Calculer $AE$ et $DE$.
\item $AD=6$, $AB=9$, $AE=10$, $AC=15$. Les droites $(DE)$ et $(BC)$ sont-elles
parallèles ? Justifier.
\item Pour mesurer la hauteur d'un arbre à Ebolowa, un élève plante un bâton vertical de
$1{,}2$ m dont l'ombre mesure $1{,}8$ m. Au même moment, l'ombre de l'arbre mesure $10{,}5$
m. Calculer la hauteur de l'arbre.
\item Pour estimer la largeur d'une rivière, un géomètre relève, dans une configuration de
Thalès, $AD=6$ m, $AB=15$ m, $AC=20$ m et $DE=8$ m, avec $(DE)\parallel(BC)$. Calculer $BC$
(largeur de la rivière) et $AE$.
\end{enumerate}

\sujetbac[papillon, réciproque et plan de maison]
\begin{enumerate}[label=\arabic*)]
\item Les droites $(AC)$ et $(BD)$ se coupent en $O$, avec $(AB)\parallel(CD)$. $OA=5$,
$OC=15$, $OB=8$. Calculer $OD$.
\item $AD=4$, $AB=9$, $AE=6$, $AC=13$. Les droites $(DE)$ et $(BC)$ sont-elles
parallèles ? Justifier.
\item Sur le plan d'une maison, une pièce a une aire de $12\text{ cm}^2$. Le plan est à
l'échelle $1/100$ (coefficient d'agrandissement réel $100$). Calculer l'aire réelle de la
pièce, en $\text{m}^2$.
\item Un jardin carré de $25\text{ m}^2$ est représenté sur un plan à l'échelle $1/50$.
Calculer l'aire du jardin sur le plan, en $\text{cm}^2$.
\end{enumerate}

% ═══════════════════════════════════════════════════════════════
\section{Corrigés}
\begin{multicols}{2}

\corrige{de l'exercice 1}
Rapport $=\dfrac69=\dfrac23$. $AE=\dfrac23\times15=10$ ; $DE=\dfrac23\times12=8$.

\corrige{de l'exercice 2}
Rapport $=\dfrac34$. $AE=\dfrac34\times12=9$ ; $DE=\dfrac34\times16=12$.

\corrige{de l'exercice 3}
Rapport $=\dfrac13$. $AE=\dfrac13\times9=3$ ; $DE=\dfrac13\times12=4$.

\corrige{de l'exercice 4}
Rapport $=\dfrac45$. $AE=\dfrac45\times10=8$ ; $DE=\dfrac45\times15=12$.

\corrige{de l'exercice 5}
$AD\times AC=3\times10=30$ ; $AB\times AE=5\times6=30$. Égaux : $(DE)\parallel(BC)$.

\corrige{de l'exercice 6}
$AD\times AC=4\times9=36$ ; $AB\times AE=6\times6=36$. Égaux : $(DE)\parallel(BC)$.

\corrige{de l'exercice 7}
$AD\times AC=2\times8=16$ ; $AB\times AE=5\times3=15$. Différents : $(DE)$ et $(BC)$ ne sont
pas parallèles.

\corrige{de l'exercice 8}
$AD\times AC=5\times10=50$ ; $AB\times AE=8\times6=48$. Différents : $(DE)$ et $(BC)$ ne
sont pas parallèles.

\corrige{de l'exercice 9}
$\dfrac{OA}{OC}=\dfrac{OB}{OD}$, donc $OD=\dfrac{OB\times OC}{OA}=\dfrac{6\times10}4=15$.

\corrige{de l'exercice 10}
Nouveau périmètre $=(4+6)\times2\times1{,}5=30$ cm. Aire initiale $=4\times6=24\text{ cm}^2$,
nouvelle aire $=24\times1{,}5^{2}=54\text{ cm}^2$.

\corrige{de l'exercice 11}
Nouveau périmètre $=18\times\dfrac23=12$ cm. Nouvelle aire $=36\times\left(\dfrac23\right)^{2}=36\times\dfrac49=16\text{ cm}^2$.

\corrige{de l'exercice 12}
Nouveau périmètre $=80\times1{,}25=100$ m.

\corrige{de l'exercice 13}
$\dfrac{h}{10{,}5}=\dfrac{1{,}2}{1{,}8}$, donc $h=10{,}5\times\dfrac{1{,}2}{1{,}8}=7$ m.

\corrige{de l'exercice 14}
Rapport $=\dfrac{AD}{AB}=\dfrac6{15}=\dfrac25$. $BC=DE\div\dfrac25=8\times\dfrac52=20$ m ;
$AE=\dfrac25\times20=8$ m.

\corrige{de l'exercice 15}
Aire réelle $=12\times100^{2}=120\,000\text{ cm}^2=12\text{ m}^2$.

\corrige{de l'exercice 16}
$\dfrac{OA}{OC}=\dfrac{OB}{OD}$, donc $OC=\dfrac{OA\times OD}{OB}=\dfrac{5\times24}8=15$.

\corrige{du sujet 1}
1) Rapport $=\dfrac8{12}=\dfrac23$. $AE=\dfrac23\times18=12$ ; $DE=\dfrac23\times15=10$.\\
2) $MR\times MP=5\times15=75$ ; $MN\times MS=8\times9=72$. Différents : $(RS)$ et $(NP)$
ne sont pas parallèles.\\
3a) Nouveau périmètre $=40\times1{,}2=48$ m.\\
3b) Nouvelle aire $=75\times1{,}2^{2}=75\times1{,}44=108\text{ m}^2$.

\corrige{du sujet 2}
1) Rapport $=\dfrac57$. $AE=\dfrac57\times14=10$ ; $DE=\dfrac57\times21=15$.\\
2) $AD\times AC=6\times15=90$ ; $AB\times AE=9\times10=90$. Égaux : $(DE)\parallel(BC)$.\\
3) $h=10{,}5\times\dfrac{1{,}2}{1{,}8}=7$ m.\\
4) Rapport $=\dfrac6{15}=\dfrac25$. $BC=8\times\dfrac52=20$ m ; $AE=\dfrac25\times20=8$ m.

\corrige{du sujet 3}
1) $OD=\dfrac{OB\times OC}{OA}=\dfrac{8\times15}5=24$.\\
2) $AD\times AC=4\times13=52$ ; $AB\times AE=9\times6=54$. Différents : $(DE)$ et $(BC)$ ne
sont pas parallèles.\\
3) Aire réelle $=12\times100^{2}=120\,000\text{ cm}^2=12\text{ m}^2$.\\
4) $25\text{ m}^2=250\,000\text{ cm}^2$. Aire sur le plan
$=250\,000\times\left(\dfrac1{50}\right)^{2}=100\text{ cm}^2$.

\end{multicols}
\exercicesBanner
